% =============================================================================
% Research: Spectral Bias, Frequency-Domain Features, and Dataset Curation
% for DETR-Based Surgical Tool Detection
% =============================================================================
% Standalone section for IEEE ACCESS article
% Compile with: \usepackage{amsmath,amssymb,booktabs,multirow,graphicx}
% =============================================================================

\section{Frequency-Domain Perspective on Dataset Curation}
\label{sec:fourier_research}

Training data composition fundamentally shapes what a neural network can
learn.  In this section, we establish a theoretical and empirical foundation
for incorporating Fourier-domain features into the dataset selection pipeline.
We draw on three converging research threads: (i)~the spectral bias of neural
networks, (ii)~frequency-domain enhancements in modern DETR architectures,
and (iii)~coverage-oriented coreset selection.


% =============================================================================
\subsection{Spectral Bias in Deep Neural Networks}
\label{sec:spectral_bias_theory}

\subsubsection{The Frequency Principle}

Rahaman et~al.~\cite{rahaman2019spectral} demonstrated that deep ReLU
networks exhibit a systematic \emph{spectral bias}: low-frequency components
of the target function are learned first, while high-frequency components
converge orders of magnitude more slowly.  This phenomenon, also known as
the \emph{frequency principle} (F-Principle), has been formalized through
the lens of the Neural Tangent Kernel (NTK)~\cite{jacot2018ntk}.

Let $f_{\boldsymbol{\theta}}: \mathbb{R}^d \to \mathbb{R}$ be a network with
parameters $\boldsymbol{\theta}$, trained on data
$\{(\mathbf{x}_i, y_i)\}_{i=1}^N$ via gradient descent with learning rate
$\eta$.  The NTK is defined as:

\begin{equation}
    \Theta(\mathbf{x}, \mathbf{x}')
    = \left\langle
        \frac{\partial f_{\boldsymbol{\theta}}(\mathbf{x})}
             {\partial \boldsymbol{\theta}},\;
        \frac{\partial f_{\boldsymbol{\theta}}(\mathbf{x}')}
             {\partial \boldsymbol{\theta}}
      \right\rangle
\end{equation}

In the infinite-width limit, $\Theta$ remains constant during training, and the
training dynamics become linear.  Decomposing the residual error into the
eigenbasis of $\Theta$:

\begin{equation}
    e_k(t) = e_k(0) \cdot \exp(-\eta \lambda_k t)
    \label{eq:ntk_convergence}
\end{equation}

where $\lambda_k$ is the $k$-th eigenvalue and $e_k(t)$ the projection of the
error onto the $k$-th eigenfunction at training step~$t$.  For standard MLPs,
the eigenvalues decay rapidly with frequency:

\begin{equation}
    \lambda_k \propto \|\boldsymbol{\omega}_k\|^{-\beta},
    \qquad \beta > 0
    \label{eq:eigenvalue_decay}
\end{equation}

where $\boldsymbol{\omega}_k$ is the frequency associated with eigenfunction
$k$.  This decay means that error components at frequency
$\|\boldsymbol{\omega}\|$ converge as
$\exp(-\eta \|\boldsymbol{\omega}\|^{-\beta} t)$ --- exponentially slower for
higher frequencies.

\subsubsection{Empirical Confirmation and Mitigation Strategies}

Recent work has both confirmed and attempted to address spectral bias:

\begin{itemize}
    \item \textbf{Frequency Regularization}~\cite{freqreg2025}: Empirical
          analysis of modern CNNs reveals that L2~regularization suppresses
          high-frequency energy accumulation by over $3\times$ compared to
          unregularized baselines, acting as an implicit low-pass filter on
          the learned features.

    \item \textbf{Multi-Grade Deep Learning}~\cite{mgdl2024}: Addresses
          spectral bias by composing multiple shallow neural networks (SNNs),
          where each grade captures low-frequency residuals of the previous
          grade.  Their compositions can represent high-frequency functions
          that monolithic networks cannot.

    \item \textbf{Fourier Feature Mapping}~\cite{tancik2020fourier}: Passing
          inputs through $\gamma(\mathbf{x}) = [\cos(2\pi\mathbf{B}\mathbf{x}),
          \sin(2\pi\mathbf{B}\mathbf{x})]$ with random frequencies
          $\mathbf{B}$ transforms the NTK into a stationary kernel with
          tunable bandwidth, enabling MLPs to learn high-frequency functions.
\end{itemize}

\subsubsection{Implication for Training Data}

Spectral bias has a direct consequence for dataset composition.
The convergence rate at frequency $\boldsymbol{\omega}$ depends not only on
the NTK eigenvalue $\lambda(\boldsymbol{\omega})$ but also on the
\emph{spectral energy of the training data} at that frequency.  Consider the
Fourier transform of the empirical data distribution:

\begin{equation}
    \hat{p}_{\text{data}}(\boldsymbol{\omega})
    = \frac{1}{N} \sum_{i=1}^{N}
      \hat{I}_i(\boldsymbol{\omega})
\end{equation}

If the training set underrepresents images with strong high-frequency content
(i.e., $\hat{p}_{\text{data}}(\boldsymbol{\omega}) \approx 0$ for large
$\|\boldsymbol{\omega}\|$), then even with NTK corrections, the network
receives insufficient gradient signal to learn high-frequency patterns.  This
motivates \emph{frequency-aware dataset curation}: selecting a training subset
that preserves the spectral distribution of the deployment domain.


% =============================================================================
\subsection{Frequency Domain in Modern DETR Architectures}
\label{sec:detr_frequency}

The importance of frequency-domain processing for object detection has been
recognized by a wave of recent DETR variants that explicitly incorporate
spectral operations into their architectures.
Table~\ref{tab:freq_detr_variants} summarizes these methods.

\begin{table*}[t]
\centering
\caption{Recent DETR variants with frequency-domain enhancement modules
         (2024--2025).  All methods demonstrate that preserving or amplifying
         high-frequency features improves detection accuracy, particularly for
         small objects.}
\label{tab:freq_detr_variants}
\begin{tabular}{llccp{5.5cm}}
\toprule
\textbf{Method} & \textbf{Venue} &
    \textbf{$\Delta$mAP} & \textbf{Key Metric} &
    \textbf{Frequency Mechanism} \\
\midrule
FDA-DETR~\cite{fda_detr2025}
    & PMC 2025
    & +4.0\,pp
    & Small obj.\ AP
    & 2D wavelet decomposition (LL/LH/HL/HH) with
      learnable band weights $(\alpha, \beta, \gamma, \delta)$ \\[4pt]
DFIR-DETR~\cite{dfir_detr2025}
    & arXiv 2025
    & --
    & 92.9\% mAP50
    & Spectral-domain constrained optimization preserving
      high-frequency boundary components \\[4pt]
Freq-DETR~\cite{freq_detr2025}
    & Expert Syst.\ 2025
    & +4.9\%
    & VisDrone mAP
    & Dual-branch spatial + frequency processing with
      FECM module \\[4pt]
FD-RTDETR~\cite{fd_rtdetr2025}
    & Electronics 2025
    & +1.9\,pp
    & COCO mAP
    & Frequency-Attention Intra-scale Feature Interaction
      for high/low frequency dual-path enhancement \\
\bottomrule
\end{tabular}
\end{table*}

\subsubsection{Wavelet Decomposition in FDA-DETR}

FDA-DETR~\cite{fda_detr2025} decomposes feature maps using the 2D discrete
wavelet transform:

\begin{equation}
    \psi_{j,m,n}(x,y) = 2^{j/2}\,\psi(2^j x - m,\; 2^j y - n)
\end{equation}

producing four sub-bands at each scale: $\mathbf{F}_{\text{LL}}$ (low-low,
global structure), $\mathbf{F}_{\text{LH}}$ (horizontal edges),
$\mathbf{F}_{\text{HL}}$ (vertical edges), and $\mathbf{F}_{\text{HH}}$
(diagonal details).  The enhanced feature map is:

\begin{equation}
    \mathbf{F}_{\text{enh}} = \alpha\,\mathbf{F}_{\text{LL}}
                            + \beta\,\mathbf{F}_{\text{LH}}
                            + \gamma\,\mathbf{F}_{\text{HL}}
                            + \delta\,\mathbf{F}_{\text{HH}}
\end{equation}

where the weights $\alpha, \beta, \gamma, \delta$ are learned per-scale.
Ablation studies show that progressively adding high-frequency components
improves small object AP from 54.8\% (LL only) to 57.5\% (all sub-bands),
confirming that \textbf{high-frequency boundary information is critical for
detecting small surgical instruments}.

\subsubsection{Spectral Optimization in DFIR-DETR}

DFIR-DETR~\cite{dfir_detr2025} formulates feature aggregation as a
constrained optimization in the spectral domain:

\begin{equation}
    \min_{\mathbf{F}}
    \|\mathcal{F}[\mathbf{F}] - \mathcal{F}[\mathbf{F}_{\text{target}}]\|^2
    + \lambda\,\Omega(\mathbf{F})
\end{equation}

where $\mathcal{F}[\cdot]$ denotes the 2D FFT and $\Omega$ is a regularizer
that penalizes loss of high-frequency components.  This directly preserves
boundary details that spatial convolutions progressively smooth away.

\subsubsection{Relevance to Dataset Curation}

These architectural innovations reveal a fundamental property: \textbf{DETR's
detection quality is bounded by the high-frequency content available in its
feature maps}.  While FDA-DETR and DFIR-DETR address this at the architecture
level, an orthogonal and complementary approach is to ensure that the
\emph{training data itself} contains sufficient high-frequency diversity.
A spectrally impoverished training set limits the features that any
architecture can extract, regardless of its frequency-enhancement modules.


% =============================================================================
\subsection{Coverage-Oriented Coreset Selection}
\label{sec:coreset_coverage}

Coreset selection aims to find a small subset $\mathcal{S} \subset \mathcal{D}$
that preserves the learning properties of the full dataset.  Recent surveys
and methods~\cite{coreset_survey2025} identify two complementary selection
criteria:

\begin{enumerate}
    \item \textbf{Difficulty-based selection}: Prioritize samples that are
          hard for the model (high loss, high EL2N
          score~\cite{paul2021deep}).  These are the most informative for
          learning decision boundaries.

    \item \textbf{Coverage-based selection}: Ensure that the selected subset
          spans the full distribution of the data manifold.  Methods include
          $k$-Center Greedy (minimax facility
          location)~\cite{coreset_survey2025}, Herding (moment
          matching)~\cite{welling2009herding}, and facility location
          (submodular coverage)~\cite{facility_location}.
\end{enumerate}

\subsubsection{The Coverage Gap in Semantic-Only Selection}

Most coreset methods operate in a single embedding space --- typically the
penultimate layer of a pre-trained classifier.  Blind Coreset
Selection~\cite{blindcs2025} uses CLIP embeddings; ELFS~\cite{elfs2025}
uses deep clustering features.  These \emph{semantic} embeddings capture
``what is in the image'' but are largely invariant to low-level statistics
such as sharpness, noise level, and frequency content.

Consider two cataract surgery frames showing the same phacoemulsification
tip in similar positions.  A DINO or CLIP embedding would map them to
nearby points ($\|\mathbf{z}_a - \mathbf{z}_b\| \approx 0$), yet one may
be in sharp focus with high $R_{\text{HL}}$ (high/low frequency ratio) and
the other blurred with low $R_{\text{HL}}$.  Semantic-only selection may
discard the blurred variant as ``redundant,'' even though it provides
valuable training signal for handling out-of-focus frames during inference.

\subsubsection{Fourier Features as a Complementary Coverage Axis}

Adding Fourier features to the selection criterion introduces a
\emph{frequency-domain coverage axis} orthogonal to semantic coverage:

\begin{equation}
    d^2(i,j) = \underbrace{\alpha_{\text{sem}}^2\,
                \|\hat{\mathbf{x}}_i^{\text{sem}} - \hat{\mathbf{x}}_j^{\text{sem}}\|^2}
               _{\text{semantic distance (DINO)}}
             + \underbrace{\alpha_{\text{four}}^2\,
                \|\hat{\mathbf{x}}_i^{\text{four}} - \hat{\mathbf{x}}_j^{\text{four}}\|^2}
               _{\text{spectral distance (Fourier)}}
    \label{eq:combined_distance}
\end{equation}

This ensures that images differing in frequency content but not in semantic
content are assigned to \emph{different} clusters, preventing spectral
homogeneity in the selected subset.

The InfoMax coreset framework~\cite{infomax2025} provides theoretical
backing: an optimal coreset maximizes the mutual information between the
subset and the full dataset.  Since Fourier features capture information
orthogonal to semantic features (Section~\ref{sec:feature_importance}),
including them strictly increases the mutual information of the selection,
improving downstream training.


% =============================================================================
\subsection{Fourier-Domain Analysis of Surgical Video Data}
\label{sec:surgical_fourier}

Cataract surgery videos present unique spectral characteristics that motivate
frequency-aware dataset curation.

\subsubsection{Sources of Frequency Variation}

\begin{enumerate}
    \item \textbf{Instrument edges}: Phacoemulsification tips, forceps, and
          IOL injectors create sharp boundaries against soft tissue,
          producing strong high-frequency content ($R_{\text{HL}} > 15$).

    \item \textbf{Specular reflections}: The wet corneal surface generates
          localized high-frequency spikes in the Fourier spectrum, with
          energy concentrated in narrow directional bands.

    \item \textbf{Tissue texture}: The iris, lens capsule, and cortical
          material exhibit characteristic mid-frequency textures
          ($0.1 \leq r < 0.5$) that vary across surgical phases.

    \item \textbf{Focus variation}: Microscope focus changes throughout
          surgery, creating frames ranging from sharp ($\hat{H}_{\text{spec}}
          \approx 0.999$, high entropy) to blurred ($\hat{H}_{\text{spec}}
          \approx 0.984$, lower entropy).

    \item \textbf{Temporal redundancy}: At 25--30\,fps, consecutive frames
          are nearly identical in both spatial and frequency domains.
          Fourier similarity filtering ($\tau_{\text{sim}} = 0.7$) efficiently
          removes these near-duplicates before expensive feature extraction.
\end{enumerate}

\subsubsection{Empirical Spectral Diversity}

Analysis of our selected dataset ($|\mathcal{S}| = 20{,}000$ images) reveals
substantial spectral diversity (Table~\ref{tab:spectral_stats}).

\begin{table}[h]
\centering
\caption{Spectral statistics of the selected training subset.
         The high coefficient of variation (CV) in $R_{\text{HL}}$ confirms
         retention of both sharp and blurred samples.}
\label{tab:spectral_stats}
\begin{tabular}{lccccc}
\toprule
\textbf{Metric} & $\min$ & $\mu$ & $\tilde{x}$ & $\sigma$ & $\max$ \\
\midrule
$R_{\text{HL}}$
    & 4.79 & 8.47 & 6.91 & 5.35 & 28.79 \\
$\hat{H}_{\text{spec}}$
    & 0.984 & 0.989 & -- & -- & 0.999 \\
$f_c$
    & 0.416 & 0.442 & -- & -- & 0.533 \\
\bottomrule
\multicolumn{6}{l}{\small CV($R_{\text{HL}}$) $= \sigma/\mu = 63.2\%$} \\
\end{tabular}
\end{table}

The $6\times$ range in $R_{\text{HL}}$ (4.79 to 28.79) spans both
\emph{smoothness-dominated} samples (blurred frames, low detail) and
\emph{frequency-dominated} samples (sharp instrument edges, specular
reflections).  Table~\ref{tab:spectral_extremes} contrasts the two extremes.

\begin{table}[h]
\centering
\caption{Spectral signature comparison: lowest vs.\ highest $R_{\text{HL}}$
         samples in the selected dataset.}
\label{tab:spectral_extremes}
\begin{tabular}{lcc}
\toprule
& \textbf{Low $R_{\text{HL}}$} & \textbf{High $R_{\text{HL}}$} \\
& (smooth, blurred) & (sharp, detailed) \\
\midrule
$R_{\text{HL}}$             & 4.79  & 28.79 \\
$\hat{E}_{\text{low}}$      & 0.072 & 0.019 \\
$\hat{E}_{\text{mid}}$      & 0.582 & 0.439 \\
$\hat{E}_{\text{high}}$     & 0.346 & 0.542 \\
$\hat{H}_{\text{spec}}$     & 0.984 & 0.999 \\
$f_c$                       & 0.416 & 0.533 \\
Anisotropy                  & 1.353 & 1.440 \\
\bottomrule
\end{tabular}
\end{table}


% =============================================================================
\subsection{Empirical Validation}
\label{sec:fourier_validation}

\subsubsection{Feature Importance Analysis}
\label{sec:feature_importance}

To quantify the contribution of each feature group to dataset selection
quality, we trained a Random Forest regressor~\cite{breiman2001random} to
predict EL2N difficulty scores from the remaining features (DINO, Fourier,
SAM).  The Gini importance results (Table~\ref{tab:rf_importance}) reveal
that Fourier features are nearly as informative as DINO embeddings for
predicting sample difficulty.

\begin{table}[h]
\centering
\caption{Random Forest feature importance for predicting EL2N difficulty.
         Per-dimension importance is computed as
         $\text{Importance}_g / d_g$.}
\label{tab:rf_importance}
\begin{tabular}{lcccc}
\toprule
\textbf{Feature} & $d_g$ & \textbf{Importance} &
    \textbf{Per-dim} & \textbf{Weight} \\
\midrule
DINO (PCA)       & 32 & 48.4\% & 1.51\% & 0.35 \\
\textbf{Fourier} &  9 & \textbf{46.3\%} & \textbf{5.14\%} & 0.15 \\
SAM              &  1 &  5.3\% & 5.30\% & 0.20 \\
\bottomrule
\end{tabular}
\end{table}

\paragraph{Key observation.}
Fourier features exhibit $5.14 / 1.51 = 3.4\times$ higher per-dimension
importance than DINO embeddings.  This disproportionate informativeness
indicates that frequency-domain statistics capture a complementary axis of
variation --- one that DINO's high-level semantic representations do not
encode.  The correlation can be explained by the spectral bias mechanism:
images that are ``difficult'' for DETR (high EL2N) typically contain either
(a)~atypical high-frequency content (sharp reflections, unusual textures) or
(b)~atypical low-frequency patterns (out-of-focus, occluded instruments)
--- both well-captured by the 9-dimensional Fourier vector.

\subsubsection{Distribution Preservation}

A necessary condition for valid dataset selection is that the selected subset
$\mathcal{S}$ preserves the statistical properties of the full pool
$\mathcal{D}$.  We verify this using the two-sample
Kolmogorov--Smirnov~(KS) test:

\begin{equation}
    D_{\text{KS}} = \sup_x \bigl|F_{\mathcal{D}}(x) - F_{\mathcal{S}}(x)\bigr|
\end{equation}

where $F_{\mathcal{D}}$ and $F_{\mathcal{S}}$ are the empirical CDFs of a
given metric over the pool ($N = 90{,}389$) and the subset
($|\mathcal{S}| = 20{,}000$), respectively.

\begin{table}[h]
\centering
\caption{Kolmogorov--Smirnov test for Fourier feature preservation.
         All $p$-values $> 0.05$: the selected subset is statistically
         indistinguishable from the full pool in the frequency domain
         despite a $4.5\times$ size reduction.}
\label{tab:ks_fourier}
\begin{tabular}{lccc}
\toprule
\textbf{Fourier Metric} &
    $\mu_{\mathcal{D}}$ & $\mu_{\mathcal{S}}$ & $p$-value \\
\midrule
$\hat{E}_{\text{low}}$      & 0.0525 & 0.0524 & 0.506 \\
$\hat{E}_{\text{mid}}$      & 0.5704 & 0.5706 & 0.820 \\
$\hat{E}_{\text{high}}$     & 0.3769 & 0.3770 & 0.928 \\
$\hat{H}_{\text{spec}}$     & 0.9894 & 0.9895 & 0.410 \\
$f_c$                       & 0.4415 & 0.4416 & 0.356 \\
$R_{\text{HL}}$             & 8.468  & 8.471  & 0.991 \\
\bottomrule
\end{tabular}
\end{table}

All $p$-values substantially exceed the $\alpha = 0.05$ significance level,
confirming that the selection procedure maintains spectral fidelity despite
reducing the dataset to 22.1\% of its original size.


% =============================================================================
\subsection{Mechanism of Action: How Fourier Features Improve DETR Training}
\label{sec:mechanism}

Integrating Fourier features into the selection criterion impacts DETR
training through three complementary mechanisms.

\subsubsection{Spectral Coverage}

By selecting samples spanning the full $R_{\text{HL}}$ range (4.79--28.79),
the training set exposes DETR to both low-frequency backgrounds (smooth
tissue) and high-frequency foregrounds (sharp tool edges).  Per
Eq.~\ref{eq:ntk_convergence}, the convergence rate for high-frequency error
components is:

\begin{equation}
    e_{\text{HF}}(t) \propto \exp\!\bigl(-\eta \lambda_{\text{HF}} t\bigr)
    \cdot \underbrace{\hat{p}_{\text{data}}(\boldsymbol{\omega}_{\text{HF}})}
          _{\text{data spectral energy}}
\end{equation}

A spectrally diverse training set increases
$\hat{p}_{\text{data}}(\boldsymbol{\omega}_{\text{HF}})$, accelerating
convergence of high-frequency features critical for boundary localization.
This is corroborated by FDA-DETR's ablation~\cite{fda_detr2025}: adding
high-frequency wavelet components (LH, HL, HH) improves small object AP by
+2.7\,pp, demonstrating the direct value of high-frequency information for
detection.

\subsubsection{Redundancy Elimination}

Surgical video at 25--30\,fps generates massive temporal redundancy.
Fourier similarity pre-filtering removes 10--20\% of near-duplicate frames
\emph{before} expensive GPU-based feature extraction (DINO, EL2N):

\begin{equation}
    \text{sim}(i,j) = \frac{1}{1 + \|\hat{\mathbf{x}}_i^{\text{four}}
                      - \hat{\mathbf{x}}_j^{\text{four}}\|_2}
\end{equation}

This serves dual purposes: (a)~reducing computational cost by $10$--$20\%$,
and (b)~preventing the $k$-means clustering from being dominated by
redundant viewing conditions, which would reduce the effective diversity of
the selected subset.

\subsubsection{Difficulty-Aware Spectral Balancing}

The high correlation between Fourier features and EL2N difficulty
(importance 46.3\%, Table~\ref{tab:rf_importance}) creates a natural
synergy: images with unusual spectral signatures tend to be more difficult
for DETR, and vice versa.  In the weighted $k$-means formulation
(Eq.~\ref{eq:combined_distance}), these spectrally unusual samples occupy
distinct regions of the feature space, forming their own clusters.  Since
one representative is selected per cluster, difficult samples with rare
spectral profiles are naturally prioritized without explicit difficulty-based
weighting.


% =============================================================================
\subsection{Connections to Related Work}
\label{sec:fourier_connections}

Table~\ref{tab:related_work} positions our approach within the broader
landscape of frequency-aware methods for training data curation and object
detection.

\begin{table*}[t]
\centering
\caption{Our approach in context: frequency-domain methods for dataset
         curation and object detection.}
\label{tab:related_work}
\begin{tabular}{p{3.2cm}p{2.5cm}p{3cm}p{5.5cm}}
\toprule
\textbf{Category} & \textbf{Representative Work} & \textbf{Venue} &
    \textbf{Relation to Our Approach} \\
\midrule
\multicolumn{4}{l}{\textit{Spectral Bias Theory}} \\
\quad Frequency Principle
    & Rahaman et~al.
    & ICML 2019
    & Motivates frequency-diverse training sets to counteract inherent
      low-frequency learning bias \\[3pt]
\quad Fourier Features
    & Tancik et~al.
    & NeurIPS 2020
    & Maps inputs to Fourier basis to tune NTK bandwidth; we apply the
      same principle to \emph{data selection} rather than architecture \\[3pt]
\quad Freq.\ Regularization
    & Zheng et~al.
    & arXiv 2025
    & Shows L2 regularization suppresses HF by $3\times$; reinforces
      need for HF-rich training data \\[3pt]
\midrule
\multicolumn{4}{l}{\textit{Frequency-Enhanced Detection}} \\
\quad FDA-DETR
    & Li et~al.
    & PMC 2025
    & Wavelet enhancement yields +4.0\,pp small obj.\ AP; confirms HF
      is the bottleneck, addressable via data or architecture \\[3pt]
\quad DFIR-DETR
    & Wang et~al.
    & arXiv 2025
    & Spectral-domain optimization preserves boundaries; we preserve
      boundaries by \emph{selecting HF-diverse training samples} \\[3pt]
\midrule
\multicolumn{4}{l}{\textit{Coreset Selection}} \\
\quad Blind CS
    & Anonymous
    & ICLR 2025
    & Embedding-space coverage for unlabeled data; we extend coverage
      to include frequency-domain embeddings \\[3pt]
\quad Coreset for OD
    & Lee et~al.
    & CVPR 2024W
    & First coreset method for object detection; uses detection-specific
      features but not frequency features \\[3pt]
\quad EL2N / Data Diet
    & Paul et~al.
    & NeurIPS 2021
    & Difficulty-based pruning; we combine with frequency-coverage for
      a balanced selection \\[3pt]
\midrule
\multicolumn{4}{l}{\textit{Medical Imaging \& Surgical Data}} \\
\quad Cataract-1K
    & Ghamsarian et~al.
    & Nat.\ Sci.\ Data 2024
    & Largest cataract surgery dataset (1000 videos); we extend with
      frequency-aware subset selection \\[3pt]
\quad Fourier Augmentation
    & Xu et~al.
    & Pattern Recogn.\ 2023
    & Amplitude-spectrum modification for domain generalization;
      confirms phase = semantics, amplitude = low-level statistics \\
\bottomrule
\end{tabular}
\end{table*}

\paragraph{Novelty of our approach.}
While frequency-domain processing has been explored both at the architecture
level (FDA-DETR, DFIR-DETR) and at the augmentation level (Fourier-based
augmentation), applying Fourier-domain features for \emph{dataset curation} in
object detection remains unexplored.  Our work bridges this gap by:
\begin{enumerate}
    \item Extracting a compact 9-dimensional Fourier signature per image that
          captures spectral diversity, directionality, and entropy;
    \item Using these features both for redundancy filtering and for
          frequency-aware clustering alongside semantic (DINO) and
          difficulty-based (EL2N) features;
    \item Demonstrating empirically that Fourier features have $3.4\times$
          higher per-dimension importance than DINO for predicting sample
          difficulty, while the selected subset preserves the spectral
          distribution of the full dataset (KS test, all $p > 0.35$).
\end{enumerate}


% =============================================================================
% BIBLIOGRAPHY ENTRIES
% =============================================================================
%
% @inproceedings{rahaman2019spectral,
%   title     = {On the Spectral Bias of Neural Networks},
%   author    = {Rahaman, Nasim and Barber, Aristide and Klartag, Boris and
%                Arjovsky, Martin and Bengio, Yoshua},
%   booktitle = {ICML},
%   year      = {2019}
% }
%
% @inproceedings{jacot2018ntk,
%   title     = {Neural Tangent Kernel: Convergence and Generalization
%                in Neural Networks},
%   author    = {Jacot, Arthur and Gabriel, Franck and Hongler, Cl{\'e}ment},
%   booktitle = {NeurIPS},
%   year      = {2018}
% }
%
% @inproceedings{tancik2020fourier,
%   title     = {Fourier Features Let Networks Learn High Frequency Functions
%                in Low Dimensional Domains},
%   author    = {Tancik, Matthew and Srinivasan, Pratul P. and Mildenhall, Ben
%                and Fridovich-Keil, Sara and Raghavan, Nithin and
%                Singhal, Utkarsh and Ramamoorthi, Ravi and Barron, Jonathan T.
%                and Ng, Ren},
%   booktitle = {NeurIPS},
%   year      = {2020}
% }
%
% @article{freqreg2025,
%   title   = {Frequency Regularization: Unveiling the Spectral Inductive Bias
%              of Deep Neural Networks},
%   author  = {Zheng, Yifei and others},
%   journal = {arXiv:2512.22192},
%   year    = {2025}
% }
%
% @inproceedings{mgdl2024,
%   title     = {Addressing Spectral Bias of Deep Neural Networks by
%                Multi-Grade Deep Learning},
%   author    = {Various},
%   booktitle = {NeurIPS},
%   year      = {2024}
% }
%
% @inproceedings{paul2021deep,
%   title     = {Deep Learning on a Data Diet: Finding Important Examples
%                Early in Training},
%   author    = {Paul, Mansheej and Ganguli, Surya and Dziugaite, Gintare K.},
%   booktitle = {NeurIPS},
%   year      = {2021}
% }
%
% @article{breiman2001random,
%   title   = {Random Forests},
%   author  = {Breiman, Leo},
%   journal = {Machine Learning},
%   volume  = {45},
%   pages   = {5--32},
%   year    = {2001}
% }
%
% @inproceedings{blindcs2025,
%   title     = {Blind Coreset Selection: Efficient Pruning for Unlabeled Data},
%   booktitle = {ICLR},
%   year      = {2025}
% }
%
% @inproceedings{elfs2025,
%   title     = {{ELFS}: Label-Free Coreset Selection with Proxy Training
%                Dynamics},
%   booktitle = {ICLR},
%   year      = {2025}
% }
%
% @inproceedings{infomax2025,
%   title     = {{InfoMax} Coreset Selection},
%   booktitle = {ICLR},
%   year      = {2025}
% }
%
% @article{coreset_survey2025,
%   title   = {A Coreset Selection of Coreset Selection Literature:
%              Introduction and Recent Advances},
%   author  = {Various},
%   journal = {arXiv:2505.17799},
%   year    = {2025}
% }
%
% @article{fda_detr2025,
%   title   = {{FDA-DETR}: A Frequency-Aware {DETR} with Dynamic Query and
%              Adaptive Multi-Task Optimization for Oriented Small Object
%              Detection},
%   author  = {Li, H. and others},
%   journal = {PMC},
%   year    = {2025}
% }
%
% @article{dfir_detr2025,
%   title   = {{DFIR-DETR}: Frequency Domain Enhancement and Dynamic Feature
%              Aggregation for Cross-Scene Small Object Detection},
%   author  = {Wang, Y. and others},
%   journal = {arXiv:2512.07078},
%   year    = {2025}
% }
%
% @article{freq_detr2025,
%   title   = {{Freq-DETR}: Frequency-Aware Transformer for Real-Time Small
%              Object Detection in Unmanned Aerial Vehicle Imagery},
%   journal = {Expert Systems with Applications},
%   year    = {2025}
% }
%
% @article{fd_rtdetr2025,
%   title   = {{FD-RTDETR}: Frequency Enhancement and Dynamic Sequence-Feature
%              Optimization for Object Detection},
%   journal = {Electronics},
%   volume  = {14},
%   number  = {23},
%   pages   = {4715},
%   year    = {2025}
% }
%
% @article{cataract1k2024,
%   title   = {Cataract-1K Dataset for Deep-Learning-Assisted Analysis of
%              Cataract Surgery Videos},
%   author  = {Ghamsarian, Negin and others},
%   journal = {Nature Scientific Data},
%   year    = {2024}
% }
%
% @article{fourier_augmentation2023,
%   title   = {Fourier-Based Augmentation with Applications to Domain
%              Generalization},
%   author  = {Xu, Q. and others},
%   journal = {Pattern Recognition},
%   year    = {2023}
% }
%
% @inproceedings{dzanic2020fourier,
%   title     = {Fourier Spectrum Discrepancies in Deep Network Generated
%                Images},
%   author    = {Dzanic, Tarik and Shah, Karan and Witherden, Freddie},
%   booktitle = {NeurIPS},
%   year      = {2020}
% }
%
% @inproceedings{coreset_od2024,
%   title     = {Coreset Selection for Object Detection},
%   author    = {Lee, Hojun and Kim, Suyoung and Lee, Junhoo},
%   booktitle = {CVPR Workshop (DDCV)},
%   year      = {2024}
% }
%
% @inproceedings{welling2009herding,
%   title     = {Herding Dynamical Weights to Learn},
%   author    = {Welling, Max},
%   booktitle = {ICML},
%   year      = {2009}
% }
